
engendrant la filtration canonique.

Ecrivant \(b=\sum a_i m_i\), avec \(a_i\in A\), et les \(m_i\) étant des monômes par rapport aux \(\gamma^j X\), en utilisant l'additivité de \(\gamma_t\), on voit qu'il suffit de faire la démonstration pour \(\gamma^k(am_i)\), avec \(a\in A\) et \(m_i\in M_i\).

En utilisant le lemme 4.11. et le \(A\)-\(\lambda\)-homomorphisme de \(B\) dans \(B\) qui envoie \(X\) sur \(m_i\), on voit qu'il suffit de faire la démonstration pour un élément du type \(\gamma^k(m_i)\).

Enfin, le polynôme \(P'_k\) étant sans composantes isobares de poids \(<k\), on est ramené au cas \(\gamma^k(\gamma^j X)\), qui résulte lui-même du fait que les polynômes \(Q'_{j,k}\) sont sans composantes isobares de poids \(<k\) (3.6).

\textbf{Remarques 4.13.} i) La décomposition de \(\Fil^k B\) donnée en 4.12 est une décomposition en somme directe de \(A\)-modules. En d'autres termes, si \(b\in\Fil^k B\), et si on écrit \(b=\sum_{j\le k}a_{i,j}m_{i,j}+b'\), avec \(a_{i,j}\in A\), \(m_{i,j}\in M_j\), et \(b'\in B\mathbin{.}M_{k+1}\), alors \(a_{i,j}\in\Fil^{k-j}A\).

ii) L'énoncé 4.12. reste vrai pour tout quotient de \(B\) par un \(\lambda\)-idéal contenu dans l'idéal d'augmentation.

iii) L'énoncé analogue à 4.12. obtenu en remplaçant l'indéterminée \(X\) par une famille d'indéterminées \((X_i)\) est encore vrai.

\textbf{Corollaire 4.14.} \(\Gr^k B\approx\Gr^k A\oplus\Gr^{k-1}A\mathbin{.}M_1\oplus\ldots\oplus\Gr^0 A\mathbin{.}M_k\).

On définit une flèche de la droite vers la gauche en associant à tout élément \(\bar a\mathbin{.}m\), où \(\bar a\) est la classe dans \(\Gr^{k-i}A\) d'un élément \(a\) de \(\Fil^{k-i}A\) et \(m\in M_i\), la classe dans \(\Gr^k B\) de \(am\). On définit une flèche en sens inverse par décomposition d'après la proposition 4.12 et la remarque 4.13.i). Il est clair que ces deux flèches sont inverses l'une de l'autre.

